% Redis-based implementation report for the project
% Generated by assistant — edit as needed
\documentclass[a4paper,11pt]{article}
\usepackage[utf8]{inputenc}
\usepackage[brazil]{babel}
\usepackage{amsmath,amssymb}
\usepackage{hyperref}
\usepackage{geometry}
\geometry{margin=1in}
\usepackage{tikz}
\usetikzlibrary{arrows.meta,positioning,shapes.geometric}
\usepackage{listings}
\usepackage{color}
\usepackage{caption}
\usepackage{enumitem}
\definecolor{codegray}{rgb}{0.95,0.95,0.95}
\lstset{backgroundcolor=\color{codegray},frame=single,breaklines=true}

\title{Implementação baseada em Redis — relatório técnico}
\author{Projeto: sistemas-distribuidos-2025.2}
\date{\today}

\begin{document}
\maketitle

\section*{Resumo}
Este documento descreve a estratégia, arquitetura, metodologia de testes (protótipos), infraestrutura e cenários no modelo SBC para a implementação baseada em Redis usada no projeto localizado na pasta \texttt{backend/} da workspace. O documento está escrito em português e foi adaptado ao código do repositório (uso de Redis para gerenciamento de sessões).

\section{Estratégia}
Objetivo: utilizar Redis como uma camada leve de armazenamento de estado (sessões) e cache para permitir sessões escaláveis e rápidas sem persistir estado sensível no servidor principal.

Principais decisões:
\begin{itemize}
  \item Usar Redis para sessões (via \texttt{connect-redis}) — sessão centralizada, compartilhada entre instâncias do backend.
  \item Executar Redis em container (Docker) para facilitar reprodução do ambiente.
  \item Proteger acesso ao Redis quando o porto do container é publicado no host (usar autenticação com \texttt{--requirepass} e publicar apenas em \texttt{127.0.0.1} durante desenvolvimento).
  \item Configurar o backend para ler a URL do Redis via variável de ambiente \texttt{REDIS_URL} (padrão: \texttt{redis://:PASSWORD@127.0.0.1:6379}).
\end{itemize}

\section{Arquitetura}
A arquitetura adotada é simples e eficaz para o propósito do projeto: Frontend (browser) \(\to\) Backend (Express) \(\to\) Redis (sessões). Abaixo está um diagrama simplificado.

\begin{figure}[h!]
\centering
\begin{tikzpicture}[node distance=2cm, every node/.style={font=\small}]
  % nodes
  \node[draw, rectangle, rounded corners, minimum width=3.5cm, minimum height=1cm] (frontend) {Frontend (React)\\http://localhost:3000};
  \node[draw, rectangle, rounded corners, minimum width=3.5cm, minimum height=1cm, right=3.5cm of frontend] (backend) {Backend (Express)\\http://localhost:3001};
  \node[draw, rectangle, rounded corners, minimum width=3.5cm, minimum height=1cm, below=1.5cm of backend] (redis) {Redis (Container)\\127.0.0.1:6379};

  % arrows
  \draw[->, thick] (frontend.east) -- node[midway,above]{XHR / fetch (axios, withCredentials)} (backend.west);
  \draw[->, thick] (backend.south) -- node[midway,right]{redis client (createClient)} (redis.north);
  \draw[->, thick] (redis.north) -- node[midway,left]{session store (connect-redis)} (backend.south);
\end{tikzpicture}
\caption{Visão arquitetural simplificada}
\end{figure}

Observações específicas do projeto:
\begin{itemize}
\item O backend atual define a URL do redis em \texttt{server.js} (exemplo: \texttt{redis://:your-strong-dev-pass@localhost:6379}). Verifique a variável de ambiente \texttt{REDIS_URL} em produção/testes.
\item As sessões são salvas com prefixo \texttt{sess:} por padrão (connect-redis). Use comandos \texttt{SCAN} / \texttt{GET} no redis-cli para inspecionar.
\end{itemize}

\section{Infraestrutura}
Recomenda-se executar Redis via Docker Compose durante o desenvolvimento. Exemplo mínimo (adaptado do repositório):

\begin{lstlisting}
# backend/docker-compose.yml (trecho)
version: '3.8'
services:
  redis:
    image: redis:7-alpine
    command:
      - redis-server
      - --bind
      - 0.0.0.0
      - --requirepass
      - "<your-dev-pass>"
      - --protected-mode
      - "yes"
    ports:
      - "127.0.0.1:6379:6379"  # bind to host loopback only for development
\end{lstlisting}

Notas de segurança:
\begin{itemize}
  \item Em desenvolvimento, publicar em \texttt{127.0.0.1} impede acesso direto da rede externa.
  \item Em produção, evite publicar Redis ao público; prefira networking interno (sem \texttt{ports:}) e autenticação forte.
  \item Habilite ACLs/TLS em ambientes expostos.
\end{itemize}

\section{Metodologia de testes (protótipos)}
Objetivo: validar comportamento de sessões, latência e tolerância a falhas para o uso de Redis. A estratégia de testes segue três níveis:

\subsection*{1. Testes unitários}
- Testes do código do backend (rotas, validação) com mocks do Redis. Use bibliotecas de testes do Node.js (Jest, Mocha).
- Mock do client Redis para verificar que a função de criação de sessão chama o método \texttt{req.session.save} corretamente.

\subsection*{2. Testes de integração (local)}
- Subir Redis via docker-compose e executar testes que fazem requests reais ao backend.
- Verificar que, após login, existe uma chave \texttt{sess:*} no Redis. Comandos úteis:
\begin{lstlisting}
# check keyspace
redis-cli -h 127.0.0.1 -p 6379 -a '<your-dev-pass>' INFO keyspace
# scan sessions
redis-cli -h 127.0.0.1 -p 6379 -a '<your-dev-pass>' --scan --pattern 'sess:*'
\end{lstlisting}

\subsection*{3. Testes de carga / protótipo}
Cenários:
\begin{enumerate}[label=\arabic*)]
  \item Conexões concorrentes de login: medir latência e throughput quando N usuários tentam logar simultaneamente (por exemplo, N=100, 500, 1000).
  \item Persistência de sessão: conectar, confirmar sessão, reiniciar backend e confirmar que a sessão persiste em Redis (se esperado).
  \item Falha e recuperação do Redis: simular reinício do container Redis e verificar comportamento do backend (reconexão, erros de sessão).
\end{enumerate}

Ferramentas de teste:
\begin{itemize}
  \item Artillery (node) ou hey / wrk para carga HTTP.
  \item redis-benchmark para testes específicos de throughput em Redis.
\end{itemize}

Exemplo simples de comando Artillery (fluxo de login):
\begin{lstlisting}
# Exemplo (instale artillery globalmente: npm i -g artillery)
# arquivo simple-login.yml
config:
  target: "http://localhost:3001"
  phases:
    - duration: 60
      arrivalRate: 20
scenarios:
  - flow:
      - post:
          url: "/api/auth/login"
          json:
            username: "user"
            password: "user123"
\end{lstlisting}

\section{Cenários no modelo SBC (interpretação)}
\textbf{Nota:} aqui adotamos uma interpretação prática de "SBC" para este documento: "Servidor (Backend) — Buffer (Redis) — Cliente (Frontend)". Ou seja, SBC descreve o fluxo entre cliente, buffer de sessão e servidor que processa lógica e autenticação.

\subsection*{Cenário 1: Login simples}
\begin{enumerate}
  \item Cliente envia POST /api/auth/login.
  \item Backend valida credenciais consultando \texttt{data/users.json} (arquivo local no projeto).
  \item Backend cria sessão em Redis (chave \texttt{sess:<id>}) usando \texttt{req.session} e responde com Set-Cookie.
  \item Cliente armazena cookie; em requests subsequentes envia cookie e backend autentica via sessão lida do Redis.
\end{enumerate}

\subsection*{Cenário 2: Multicliente concorrente}
\begin{itemize}
  \item Objetivo: observar impacto no tempo de resposta e latência por sessão quando 100/500 clientes simultâneos fazem login.
  \item Métricas: throughput, latência 50/95/99 perc., taxas de erro, número de chaves no Redis.
\end{itemize}

\subsection*{Cenário 3: Reinício do backend (sem perder sessão)}
\begin{enumerate}
  \item Com sessão criada no Redis, reinicie o processo do backend.
  \item Verificar que clientes com cookie válido continuam autenticados (backend lê a sessão do Redis).
\end{enumerate}

\subsection*{Cenário 4: Falha temporária do Redis}
\begin{enumerate}
  \item Simule parada do Redis (docker stop) e observe as reações do backend (logs, erros de conexão).
  \item Verificar que operações críticas falham com mensagens de erro apropriadas e que, após o reinício do Redis, o backend reconecta e restaura funcionalidade.
\end{enumerate}

\section{Checklist de implantação / comandos úteis}
\begin{enumerate}
  \item Subir Redis no host (bind em loopback e com senha para dev):
  \begin{lstlisting}
  sudo docker-compose -f backend/docker-compose.yml up -d
  \end{lstlisting}
  \item Rodar backend no host (usa REDIS_URL):
  \begin{lstlisting}
  # no diretório backend
  export REDIS_URL='redis://:your-strong-dev-pass@127.0.0.1:6379'
  npm start
  \end{lstlisting}
  \item Testar login via curl:
  \begin{lstlisting}
  curl -c cookies.txt -H "Content-Type: application/json" -d '{"username":"admin","password":"admin123"}' http://localhost:3001/api/auth/login
  \end{lstlisting}
\end{enumerate}

\section*{Conclusão e próximos passos}
\begin{itemize}
  \item Esta proposta oferece um roteiro para usar Redis como buffer de sessões no projeto com segurança local (bind em 127.0.0.1) e testes práticos para validar comportamento sob carga.
  \item Próximo passo opcional: eu posso gerar um arquivo de configuração Artillery para os cenários de carga e patchar \texttt{backend/server.js} para usar \texttt{process.env.REDIS_URL} em produção.
\end{itemize}

\vfill
\noindent Documento gerado automaticamente — revise as senhas e variáveis antes de publicar.

\end{document}
